% ── SECTION 5: QUANTUM PHYSICS AS CONFIRMATION, NOT APPROPRIATION (~500 words)

\section{Quantum Physics as Confirmation, Not Appropriation}
\label{sec:confirmation}

% This project does not claim \Bahaullah was writing physics.
% Claim: the structural description in the writings CORRESPONDS to what
% physics discovered independently.
% Relationship: confirmation, not borrowing.
% Each tradition keeps its own vocabulary and methodology.

A clarification is essential.

This paper does not claim that \Bahaullah was writing physics, or that the
\bahai writings anticipate the technical formalism of quantum mechanics.
The writings are not a physics textbook.
They are spiritual revelation, addressed to questions of human meaning,
moral development, and the nature of God.

What this paper claims is different: that the \emph{structural description}
of reality in the \bahai writings --- existence as relational affinity, creation
as relational event, love as the creative foundation at every degree --- corresponds
to the structural description that quantum mechanics has independently arrived at
through a century of mathematics and experiment.

The relationship is confirmation, not appropriation.
The \bahai writings do not need quantum mechanics to be true.
Quantum mechanics does not need the \bahai writings to be true.
What we document is that, independently of each other, both arrived at
the same description --- which is what the two-wings principle predicts they should do.

Each tradition keeps its own vocabulary and methodology.
A \bahai scholar does not need to accept the formalism of the Relational Coherence
Theorem to validate this paper's argument.
A physicist does not need to accept the authority of the \bahai writings.
Each tradition validates the convergence on its own grounds.
